\documentclass{article}
\usepackage[margin=1.5cm, includefoot, footskip=30pt]{geometry}
\usepackage{amsmath}
\usepackage{amsfonts}
\usepackage{amssymb}
\usepackage{graphicx}
\usepackage{hyperref}
\usepackage{subcaption}
\usepackage{tikz}
\usetikzlibrary{calc}
\usetikzlibrary{shapes.geometric}
\usetikzlibrary{arrows.meta, positioning}
\usepackage{latexsym}
\newcommand{\lt}{\ensuremath <}
\newcommand{\gt}{\ensuremath >}
\usepackage{float}
\usepackage{framed}
\usetikzlibrary{positioning}
\usetikzlibrary{matrix}
\usepackage{amsthm}

\title{Supplementary Information}
\author{Harry Foster, Vince Knight, Sebastian Krapohl}
\date{\today}

\begin{document}
\maketitle

\section{Parameter Sweep}
\begin{figure}[htbp]
    \centering
    \begin{tabular}{|c|c|c|c|}
        \hline
        Parameter & Minimum & Maximum & Number of Values per $N$\\
        \hline
        $N$ & 3 & 8 & N/A\\
        \hline
        $r$ & $\frac{1}{2}$ & $\frac{3N}{2}$ & 10\\
        \hline
        $M$ & $N$ & $4N$ & 10\\
        \hline
        $\mu$ & 0.001 & 0.1 & 4\\
        \hline
        $\epsilon$ & 0 & $0.99 \cdot \max_i \alpha_i$ &10\\
        \hline
        $\beta$ & 0 & 2 & 10\\
        \hline
        $A$ & 1 & \(\frac{M}{N}\) & 10\\
        \hline
    \end{tabular}
    \caption{The set of measured parameters and boundry values. Some parameters
    are only applicable to certain processes. The value $A$ is the
    \emph{aspiration} considered in aspiration dynamics, in section~4}
    \label{fig:parameter_sweep_table}
\end{figure}


\section{Explicit transition matrices for the two-player game}
\label{sec:explicit_matrices}

In each subsection below, states are ordered
\(\mathbf{a} = (0,0),\ \mathbf{b} = (0,1),\ \mathbf{c} = (1,0),\
\mathbf{d} = (1,1)\); rows and columns of each matrix correspond to these
states in the same order.

\subsection*{Moran process}

\begin{equation}
    \resizebox{\textwidth}{!}{$
    \begin{bmatrix}
        1 & 0 & 0 & 0\\
        \frac{f_{1}(\mathbf{b})}{2 f_{1}(\mathbf{b}) + 2 f_{2}(\mathbf{b})} &
        1 - \frac{f_{1}(\mathbf{b})}{2 f_{1}(\mathbf{b}) + 2 f_{2}(\mathbf{b})}
        - \frac{f_{2}(\mathbf{b})}{2 f_{1}(\mathbf{b}) + 2 f_{2}(\mathbf{b})} &
        0 &
        \frac{f_{2}(\mathbf{b})}{2 f_{1}(\mathbf{b}) + 2 f_{2}(\mathbf{b})}\\
        \frac{f_{2}(\mathbf{c})}{2 f_{1}(\mathbf{c}) + 2 f_{2}(\mathbf{c})} &
        0 &
        1 - \frac{f_{1}(\mathbf{c})}{2 f_{1}(\mathbf{c}) + 2 f_{2}(\mathbf{c})}
        - \frac{f_{2}(\mathbf{c})}{2 f_{1}(\mathbf{c}) + 2 f_{2}(\mathbf{c})} &
        \frac{f_{1}(\mathbf{c})}{2 f_{1}(\mathbf{c}) + 2 f_{2}(\mathbf{c})}\\
        0 & 0 & 0 & 1
    \end{bmatrix}$}
    \label{eqn:2x2_Moran_Transition}
\end{equation}


Now, by taking the fundamental matrix of \(T\) and performing a right
multiplication by \(R\), the matrix with entries which represent transitions
from transitive states to absorbing
states\cite{f66ea2a5-f1c3-3855-8dc0-e3047d83298e}, we can obtain the absorption
matrix for this system.

For the above game, we find the following equation:

\begin{equation}
    \left(I - \begin{bmatrix}1 - \frac{f_{1}{\left(\textbf{b} \right)}}{2 f_{1}{\left(\textbf{b} \right)} + 2 f_{2}{\left(\textbf{b} \right)}} - \frac{f_{2}{\left(\textbf{b} \right)}}{2 f_{1}{\left(\textbf{b} \right)} + 2 f_{2}{\left(\textbf{b} \right)}} & 0\\0 & 1 - \frac{f_{1}{\left(\textbf{c} \right)}}{2 f_{1}{\left(\textbf{c} \right)} + 2 f_{2}{\left(\textbf{c} \right)}} - \frac{f_{2}{\left(\textbf{c} \right)}}{2 f_{1}{\left(\textbf{c} \right)} + 2 f_{2}{\left(\textbf{c} \right)}}\end{bmatrix}\right)^{-1}
        \begin{bmatrix}
            \frac{f_{1}{\left(\textbf{b} \right)}}{2 f_{1}{\left(\textbf{b} \right)} + 2 f_{2}{\left(\textbf{b} \right)}} & \frac{f_{2}{\left(\textbf{b} \right)}}{2 f_{1}{\left(\textbf{b} \right)} + 2 f_{2}{\left(\textbf{b} \right)}}\\\frac{f_{2}{\left(\textbf{c} \right)}}{2 f_{1}{\left(\textbf{c} \right)} + 2 f_{2}{\left(\textbf{c} \right)}} & \frac{f_{1}{\left(\textbf{c} \right)}}{2 f_{1}{\left(\textbf{c} \right)} + 2 f_{2}{\left(\textbf{c} \right)}}
        \end{bmatrix}
\end{equation}

And solve to find our absorption matrix:

\begin{equation}
    \begin{bmatrix}
        \frac{f_{1}{\left(\textbf{b} \right)}}{\left(\frac{f_{1}{\left(\textbf{b} \right)}}{2 f_{1}{\left(\textbf{b} \right)} + 2 f_{2}{\left(\textbf{b} \right)}} + \frac{f_{2}{\left(\textbf{b} \right)}}{2 f_{1}{\left(\textbf{b} \right)} + 2 f_{2}{\left(\textbf{b} \right)}}\right) \left(2 f_{1}{\left(\textbf{b} \right)} + 2 f_{2}{\left(\textbf{b} \right)}\right)} & \frac{f_{2}{\left(\textbf{b} \right)}}{\left(\frac{f_{1}{\left(\textbf{b} \right)}}{2 f_{1}{\left(\textbf{b} \right)} + 2 f_{2}{\left(\textbf{b} \right)}} + \frac{f_{2}{\left(\textbf{b} \right)}}{2 f_{1}{\left(\textbf{b} \right)} + 2 f_{2}{\left(\textbf{b} \right)}}\right) \left(2 f_{1}{\left(\textbf{b} \right)} + 2 f_{2}{\left(\textbf{b} \right)}\right)}\\\frac{f_{2}{\left(\textbf{c} \right)}}{\left(\frac{f_{1}{\left(\textbf{c} \right)}}{2 f_{1}{\left(\textbf{c} \right)} + 2 f_{2}{\left(\textbf{c} \right)}} + \frac{f_{2}{\left(\textbf{c} \right)}}{2 f_{1}{\left(\textbf{c} \right)} + 2 f_{2}{\left(\textbf{c} \right)}}\right) \left(2 f_{1}{\left(\textbf{c} \right)} + 2 f_{2}{\left(\textbf{c} \right)}\right)} & \frac{f_{1}{\left(\textbf{c} \right)}}{\left(\frac{f_{1}{\left(\textbf{c} \right)}}{2 f_{1}{\left(\textbf{c} \right)} + 2 f_{2}{\left(\textbf{c} \right)}} + \frac{f_{2}{\left(\textbf{c} \right)}}{2 f_{1}{\left(\textbf{c} \right)} + 2 f_{2}{\left(\textbf{c} \right)}}\right) \left(2 f_{1}{\left(\textbf{c} \right)} + 2 f_{2}{\left(\textbf{c} \right)}\right)}
    \end{bmatrix}
    \label{eqn:2x2_Moran_Absorption}
\end{equation}

In this matrix, the rows represent our initial transitive state (in our system,
row 1 represents \textbf{b} and row 2 represents \textbf{c}), and our columns
represent which state the system is absorbed into (with \textbf{a} on the left,
and \textbf{d} on the right). It is immediately clear from this result that the
probability of the system ending up in a particular absorbing state does not
rely on the fitness of individuals in the absorbing state, but rather their
fitness in the transitive states. Thus, players can remain in a state which
could be bettered by a single change, however they will not make such a change
if no individual of another type exists in the population.

\subsection*{Fermi imitation dynamics}
\begin{center}
\begin{equation}
\begin{bmatrix}
1 & 0 & 0 & 0\\
\frac{1}{2}\phi(\pi_2(\mathbf{b}) - \pi_1(\mathbf{b})) &
\frac{1}{2} &
0 &
\frac{1}{2}\phi(\pi_1(\mathbf{b}) - \pi_2(\mathbf{b}))\\
\frac{1}{2}\phi(\pi_1(\mathbf{c}) - \pi_2(\mathbf{c})) &
0 &
\frac{1}{2} &
\frac{1}{2}\phi(\pi_2(\mathbf{c}) - \pi_1(\mathbf{c}))\\
0 & 0 & 0 & 1
\end{bmatrix}
\label{eqn:2x2_Fermi_Transition}
\end{equation}
\end{center}

For the two-player state space \(\mathbf{a},\mathbf{b},\mathbf{c},\mathbf{d}\),
the explicit transition matrix is given in
equation~(\ref{eqn:2x2_Fermi_Transition}) in the supplementary material.
Applying the same absorption matrix calculation as for the Moran process, we
obtain:

\begin{equation}
    \begin{bmatrix}
        \phi(\pi_2(\textbf{b}) - \pi_1(\textbf{b})) & \phi(\pi_1(\textbf{b})-\pi_2(\textbf{b}))\\\phi(\pi_1(\textbf{c}) - \pi_2(\textbf{c})) & \phi(\pi_2(\textbf{c}) - \pi_1(\textbf{c}))
    \end{bmatrix}
\end{equation}

This shows that, similarly to the Moran process, our absorption probabilities
rely on the fitness of players in transitive states and not those of players in
the absorbing states.
Fermi imitation dynamics is therefore poorly suited to heterogeneous
populations: a player may copy a high-fitness strategy without considering
whether that strategy would benefit them, and may even enter an absorbing state
that worsens their own payoff.

\subsection*{Aspiration Dynamics}

The transition matrix for the two-player, two-action case with state space
\(\mathbf{a} = (0,0),\ \mathbf{b} = (0,1),\ \mathbf{c} = (1,0),\ \mathbf{d} = (1,1)\) is:
\begin{equation}
\resizebox{\textwidth}{!}{
$
\begin{bmatrix}
1-\frac{1}{2(1+e^{-\beta(A_2-f_2(a))})}-\frac{1}{2(1+e^{-\beta(A_1-f_1(a))})} &
\frac{e^{A_2\beta}}{2(e^{A_2\beta}+e^{\beta f_2(a)})} &
\frac{e^{A_1\beta}}{2(e^{A_1\beta}+e^{\beta f_1(a)})} &
0 \\

\frac{e^{A_2\beta}}{2(e^{A_2\beta}+e^{\beta f_2(b)})} &
1-\frac{1}{2(1+e^{-\beta(A_2-f_2(b))})}-\frac{1}{2(1+e^{-\beta(A_1-f_1(b))})} &
0 &
\frac{e^{A_1\beta}}{2(e^{A_1\beta}+e^{\beta f_1(b)})} \\

\frac{e^{A_1\beta}}{2(e^{A_1\beta}+e^{\beta f_1(c)})} &
0 &
1-\frac{1}{2(1+e^{-\beta(A_2-f_2(c))})}-\frac{1}{2(1+e^{-\beta(A_1-f_1(c))})} &
\frac{e^{A_2\beta}}{2(e^{A_2\beta}+e^{\beta f_2(c)})} \\

0 &
\frac{e^{A_1\beta}}{2(e^{A_1\beta}+e^{\beta f_1(d)})} &
\frac{e^{A_2\beta}}{2(e^{A_2\beta}+e^{\beta f_2(d)})} &
1-\frac{1}{2(1+e^{-\beta(A_2-f_2(d))})}-\frac{1}{2(1+e^{-\beta(A_1-f_1(d))})}
\end{bmatrix}
$
}
\label{eqn:2x2_Aspiration_Transition}
\end{equation}
We can see that as our players no longer consider the actions of the rest of
the population when considering their decision, we no longer have an absorbing
Markov chain.
Thus, we no longer consider an absorption matrix, but instead we consider the
steady state of our Markov chain \cite{f66ea2a5-f1c3-3855-8dc0-e3047d83298e}
This is the proportion of time spent in each state when the Markov chain
reaches a stable distribution across the states.
Any irreducible, aperiodic, finite Markov chain has a unique stationary
distribution, and the chain converges to it from any starting
state~\cite{levin2017markov}..
For our four state system, we acquire the following result for a stationary
distribution:

%let thing finish running

\subsection*{Introspective Dynamics}

The transition matrix for the two-player, two-action case with state space
\(\mathbf{a} = (0,0),\ \mathbf{b} = (0,1),\ \mathbf{c} = (1,0),\ \mathbf{d} = (1,1)\),
following Couto~\cite{Couto2022}, is:

\begin{equation}
\resizebox{\textwidth}{!}{$
\begin{bmatrix}
    1 - \frac{1}{2\left(e^{\beta\left(\pi_{2}(a) - \pi_{2}(b)\right)} + 1\right)} - \frac{1}{2\left(e^{\beta\left(\pi_{1}(a) - \pi_{1}(c)\right)} + 1\right)} 
    & \frac{1}{2\left(e^{\beta\left(\pi_{2}(a) - \pi_{2}(b)\right)} + 1\right)} 
    & \frac{1}{2\left(e^{\beta\left(\pi_{1}(a) - \pi_{1}(c)\right)} + 1\right)} 
    & 0 \\[6pt]

    \frac{1}{2\left(e^{\beta\left(- \pi_{2}(a) + \pi_{2}(b)\right)} + 1\right)} 
    & 1 - \frac{1}{2\left(e^{\beta\left(\pi_{1}(b) - \pi_{1}(d)\right)} + 1\right)} - \frac{1}{2\left(1 + e^{-\beta\left(\pi_{2}(a) - \pi_{2}(b)\right)}\right)} 
    & 0 
    & \frac{1}{2\left(e^{\beta\left(\pi_{1}(b) - \pi_{1}(d)\right)} + 1\right)} \\[6pt]

    \frac{1}{2\left(e^{\beta\left(- \pi_{1}(a) + \pi_{1}(c)\right)} + 1\right)} 
    & 0 
    & 1 - \frac{1}{2\left(e^{\beta\left(\pi_{2}(c) - \pi_{2}(d)\right)} + 1\right)} - \frac{1}{2\left(1 + e^{-\beta\left(\pi_{1}(a) - \pi_{1}(c)\right)}\right)} 
    & \frac{1}{2\left(e^{\beta\left(\pi_{2}(c) - \pi_{2}(d)\right)} + 1\right)} \\[6pt]

    0 
    & \frac{1}{2\left(e^{\beta\left(- \pi_{1}(b) + \pi_{1}(d)\right)} + 1\right)} 
    & \frac{1}{2\left(e^{\beta\left(- \pi_{2}(c) + \pi_{2}(d)\right)} + 1\right)} 
    & 1 - \frac{1}{2\left(1 + e^{-\beta\left(\pi_{2}(c) - \pi_{2}(d)\right)}\right)} - \frac{1}{2\left(1 + e^{-\beta\left(\pi_{1}(b) - \pi_{1}(d)\right)}\right)}
\end{bmatrix}$
}
\label{eqn:2x2_Introspection_Transition}
\end{equation}

We then consider the steady state vectors for this system.
For this, we take the left eigenvector of the transition matrix, which is given by:




Unlike in the case of imitation dynamics, there are no absorbing states for a Markov chain transitioning according to imitation dynamics.
Therefore, we define a steady state vector \textbf{u} according to the method Couto uses in \cite{Couto2022}. 



\subsection*{Introspective Imitation Dynamics}

Note that Introspective Imitation Dynamics takes into
account both selection intensity and choice intensity. This means that both the
level to which a player's fitness is considered when their strategy is being chosen for
consideration, and also how much it is considered when choosing whether to
accept the new strategy, is present in the dynamic.

Let \(\eta_{i,j}(\mathbf{x}, \mathbf{y}) = 1 + e^{\beta(f_i(\mathbf{x}) - f_j(\mathbf{y}))}\)
be the denominator of the Fermi function. For the two-player state space
\(S = \{\mathbf{a},\mathbf{b},\mathbf{c},\mathbf{d}\}\), the explicit
transition matrix is given in equation~(\ref{eqn:2x2_IID_Transition}) in the
supplementary material. Applying the absorption matrix calculation, we obtain:

\begin{center}
\begin{large}
\begin{equation}
\begin{bmatrix}
\frac{\eta_{1,1}(\textbf{b,d})f_{1}{\left(b \right)}}{\eta_{1,1}(\textbf{b,d}) f_{1}{\left(b \right)} + \eta_{2,2}(\textbf{b,a}) f_{2}{\left(b \right)}} & \frac{\eta_{2,2}(\textbf{b,a}) f_{2}{\left(b \right)}}{\eta_{1,1}(\textbf{b,d}) f_{1}{\left(b \right)} + \eta_{2,2}(\textbf{b,a}) f_{2}{\left(b \right)}}\\\frac{\eta_{2,2}(\textbf{c,d}) f_{2}{\left(c \right)}}{\eta_{1,1}(\textbf{c,a}) f_{1}{\left(c \right)} + \eta_{2,2}(\textbf{c,d}) f_{2}{\left(c \right)}} & \frac{\eta_{1,1}(\textbf{c,a}) f_{1}{\left(c \right)}}{\eta_{1,1}(\textbf{c,a}) f_{1}{\left(c \right)} + \eta_{2,2}(\textbf{c,d}) f_{2}{\left(c \right)}}
\end{bmatrix}
\end{equation}
\end{large}
\end{center}

Unlike other imitation dynamics, we can see that the absorption probability in an introspective imitation process relies on the fitness that at least one individual will possess in the absorbing state.
This means that players will enter absorbing states with a higher probability if it improves their own fitness, despite any heterogeneity between payoffs for the same action played by different players.
This solves the problem of players making worsening moves based on the fitness of another individual in the population, while still retaining the property of imitation dynamics that new strategies cannot be introduced to the system.

Let \(\eta_{i,j}(\mathbf{x}, \mathbf{y}) = 1 + e^{\beta(f_i(\mathbf{x}) - f_j(\mathbf{y}))}\).
\begin{center}
\begin{equation}
\resizebox{\textwidth}{!}{$
\begin{bmatrix}
1 & 0 & 0 & 0\\
\frac{f_{1}(b)}{2\eta_{2,2}(\mathbf{b},\mathbf{a})(f_{1}(b)+f_{2}(b))} &
1 - \frac{f_{2}(b)}{\eta_{1,1}(\mathbf{b},\mathbf{d})(2f_{1}(b)+2f_{2}(b))}
  - \frac{f_{1}(b)}{\eta_{2,2}(\mathbf{b},\mathbf{a})(2f_{1}(b)+2f_{2}(b))} &
0 &
\frac{f_{2}(b)}{2\eta_{1,1}(\mathbf{b},\mathbf{d})(f_{1}(b)+f_{2}(b))}\\
\frac{f_{2}(c)}{2\eta_{1,1}(\mathbf{c},\mathbf{a})(f_{1}(c)+f_{2}(c))} &
0 &
1 - \frac{f_{1}(c)}{\eta_{2,2}(\mathbf{c},\mathbf{d})(2f_{1}(c)+2f_{2}(c))}
  - \frac{f_{2}(c)}{\eta_{1,1}(\mathbf{c},\mathbf{a})(2f_{1}(c)+2f_{2}(c))} &
\frac{f_{1}(c)}{2\eta_{2,2}(\mathbf{c},\mathbf{d})(f_{1}(c)+f_{2}(c))}\\
0 & 0 & 0 & 1
\end{bmatrix}$}
\label{eqn:2x2_IID_Transition}
\end{equation}
\end{center}

\subsection*{Moran process with mutation}
\begin{equation}
\resizebox{\textwidth}{!}{$
\begin{bmatrix}
-\frac{\mu_{12}}{2} - \frac{\mu_{22}}{2} + 1 &
\frac{\mu_{22}}{2} &
\frac{\mu_{12}}{2} &
0\\[0.5em]
\frac{\mu_{21}}{2}
+ \frac{(\epsilon f_{1}(b) + 1)(-\mu_{21} - \mu_{22} + 1)}
       {2(\epsilon f_{1}(b) + \epsilon f_{2}(b) + 2)}
&
-\frac{\mu_{12}}{2} - \frac{\mu_{21}}{2}
- \frac{(\epsilon f_{1}(b) + 1)(-\mu_{21} - \mu_{22} + 1)}
       {2(\epsilon f_{1}(b) + \epsilon f_{2}(b) + 2)}
- \frac{(\epsilon f_{2}(b) + 1)(-\mu_{11} - \mu_{12} + 1)}
       {2(\epsilon f_{1}(b) + \epsilon f_{2}(b) + 2)}
+ 1 &
0 &
\frac{\mu_{12}}{2}
+ \frac{(\epsilon f_{2}(b) + 1)(-\mu_{11} - \mu_{12} + 1)}
       {2(\epsilon f_{1}(b) + \epsilon f_{2}(b) + 2)}\\[0.5em]
\frac{\mu_{11}}{2}
+ \frac{(\epsilon f_{2}(c) + 1)(-\mu_{11} - \mu_{12} + 1)}
       {2(\epsilon f_{1}(c) + \epsilon f_{2}(c) + 2)}
& 0 &
-\frac{\mu_{11}}{2} - \frac{\mu_{22}}{2}
- \frac{(\epsilon f_{1}(c) + 1)(-\mu_{21} - \mu_{22} + 1)}
       {2(\epsilon f_{1}(c) + \epsilon f_{2}(c) + 2)}
- \frac{(\epsilon f_{2}(c) + 1)(-\mu_{11} - \mu_{12} + 1)}
       {2(\epsilon f_{1}(c) + \epsilon f_{2}(c) + 2)}
+ 1 &
\frac{\mu_{22}}{2}
+ \frac{(\epsilon f_{1}(c) + 1)(-\mu_{21} - \mu_{22} + 1)}
       {2(\epsilon f_{1}(c) + \epsilon f_{2}(c) + 2)}\\[0.5em]
0 & \frac{\mu_{11}}{2} & \frac{\mu_{21}}{2} &
-\frac{\mu_{11}}{2} - \frac{\mu_{21}}{2} + 1
\end{bmatrix}$}
\label{eqn:2x2_Moran_Mutation_Transition}
\end{equation}

\subsection*{Fermi imitation dynamics with mutation}
\begin{equation}
\resizebox{\textwidth}{!}{$
\begin{bmatrix}
-\frac{\mu_{12}}{2} - \frac{\mu_{22}}{2} + 1 &
\frac{\mu_{22}}{2} & \frac{\mu_{12}}{2} & 0\\[1em]
\frac{\mu_{21}}{2}
+ \frac{-\mu_{21} - \mu_{22} + 1}{2(e^{\beta(-f_{1}(b) + f_{2}(b))} + 1)}
&
-\frac{\mu_{12}}{2} - \frac{\mu_{21}}{2} + 1
- \frac{-\mu_{11} - \mu_{12} + 1}{2(e^{\beta(f_{1}(b) - f_{2}(b))} + 1)}
- \frac{-\mu_{21} - \mu_{22} + 1}{2(e^{\beta(-f_{1}(b) + f_{2}(b))} + 1)}
& 0 &
\frac{\mu_{12}}{2}
+ \frac{-\mu_{11} - \mu_{12} + 1}{2(e^{\beta(f_{1}(b) - f_{2}(b))} + 1)}\\[1em]
\frac{\mu_{11}}{2}
+ \frac{-\mu_{11} - \mu_{12} + 1}{2(e^{\beta(f_{1}(c) - f_{2}(c))} + 1)}
& 0 &
-\frac{\mu_{11}}{2} - \frac{\mu_{22}}{2} + 1
- \frac{-\mu_{11} - \mu_{12} + 1}{2(e^{\beta(f_{1}(c) - f_{2}(c))} + 1)}
- \frac{-\mu_{21} - \mu_{22} + 1}{2(e^{\beta(-f_{1}(c) + f_{2}(c))} + 1)}
&
\frac{\mu_{22}}{2}
+ \frac{-\mu_{21} - \mu_{22} + 1}{2(e^{\beta(-f_{1}(c) + f_{2}(c))} + 1)}\\[1em]
0 & \frac{\mu_{11}}{2} & \frac{\mu_{21}}{2} &
-\frac{\mu_{11}}{2} - \frac{\mu_{21}}{2} + 1
\end{bmatrix}$}
\label{eqn:2x2_Fermi_Mutation_Transition}
\end{equation}

\subsection*{Aspiration dynamics with mutation}
\begin{equation}
\resizebox{\textwidth}{!}{$
\left[\begin{matrix}
1 - \frac{0.5}{1 + e^{-\beta(A_{2} - f_{2}(a))}} - \frac{0.5}{1 + e^{-\beta(A_{1} - f_{1}(a))}} &
\frac{0.5 e^{A_{2}\beta}}{e^{A_{2}\beta} + e^{\beta f_{2}(a)}} &
\frac{0.5 e^{A_{1}\beta}}{e^{A_{1}\beta} + e^{\beta f_{1}(a)}} & 0\\
\frac{0.5 e^{A_{2}\beta}}{e^{A_{2}\beta} + e^{\beta f_{2}(b)}} &
1 - \frac{0.5}{1 + e^{-\beta(A_{2} - f_{2}(b))}} - \frac{0.5}{1 + e^{-\beta(A_{1} - f_{1}(b))}} &
0 &
\frac{0.5 e^{A_{1}\beta}}{e^{A_{1}\beta} + e^{\beta f_{1}(b)}}\\
\frac{0.5 e^{A_{1}\beta}}{e^{A_{1}\beta} + e^{\beta f_{1}(c)}} & 0 &
1 - \frac{0.5}{1 + e^{-\beta(A_{2} - f_{2}(c))}} - \frac{0.5}{1 + e^{-\beta(A_{1} - f_{1}(c))}} &
\frac{0.5 e^{A_{2}\beta}}{e^{A_{2}\beta} + e^{\beta f_{2}(c)}}\\
0 &
\frac{0.5 e^{A_{1}\beta}}{e^{A_{1}\beta} + e^{\beta f_{1}(d)}} &
\frac{0.5 e^{A_{2}\beta}}{e^{A_{2}\beta} + e^{\beta f_{2}(d)}} &
1 - \frac{0.5}{1 + e^{-\beta(A_{2} - f_{2}(d))}} - \frac{0.5}{1 + e^{-\beta(A_{1} - f_{1}(d))}}
\end{matrix}\right]$}
\label{eqn:2x2_Aspiration_Mutation_Transition}
\end{equation}

\subsection*{Introspection dynamics with mutation}
\begin{equation}
\resizebox{\textwidth}{!}{$
\begin{bmatrix}
-\frac{\mu_{12}}{2} - \frac{\mu_{22}}{2} + 1
- \frac{-\mu_{21} - \mu_{22} + 1}{2(e^{\beta(f_{2}(a) - f_{2}(b))} + 1)}
- \frac{-\mu_{11} - \mu_{12} + 1}{2(e^{\beta(f_{1}(a) - f_{1}(c))} + 1)} &
\frac{\mu_{22}}{2} + \frac{-\mu_{21} - \mu_{22} + 1}{2(e^{\beta(f_{2}(a) - f_{2}(b))} + 1)} &
\frac{\mu_{12}}{2} + \frac{-\mu_{11} - \mu_{12} + 1}{2(e^{\beta(f_{1}(a) - f_{1}(c))} + 1)} &
0\\[1em]
\frac{\mu_{21}}{2} + \frac{-\mu_{21} - \mu_{22} + 1}{2(e^{\beta(-f_{2}(a) + f_{2}(b))} + 1)} &
-\frac{\mu_{12}}{2} - \frac{\mu_{21}}{2} + 1
- \frac{-\mu_{21} - \mu_{22} + 1}{2(e^{\beta(-f_{2}(a) + f_{2}(b))} + 1)}
- \frac{-\mu_{11} - \mu_{12} + 1}{2(e^{\beta(f_{1}(b) - f_{1}(d))} + 1)} &
0 &
\frac{\mu_{12}}{2} + \frac{-\mu_{11} - \mu_{12} + 1}{2(e^{\beta(f_{1}(b) - f_{1}(d))} + 1)}\\[1em]
\frac{\mu_{11}}{2} + \frac{-\mu_{11} - \mu_{12} + 1}{2(e^{\beta(-f_{1}(a) + f_{1}(c))} + 1)} &
0 &
-\frac{\mu_{11}}{2} - \frac{\mu_{22}}{2} + 1
- \frac{-\mu_{21} - \mu_{22} + 1}{2(e^{\beta(f_{2}(c) - f_{2}(d))} + 1)}
- \frac{-\mu_{11} - \mu_{12} + 1}{2(e^{\beta(-f_{1}(a) + f_{1}(c))} + 1)} &
\frac{\mu_{22}}{2} + \frac{-\mu_{21} - \mu_{22} + 1}{2(e^{\beta(f_{2}(c) - f_{2}(d))} + 1)}\\[1em]
0 &
\frac{\mu_{11}}{2} + \frac{-\mu_{11} - \mu_{12} + 1}{2(e^{\beta(-f_{1}(b) + f_{1}(d))} + 1)} &
\frac{\mu_{21}}{2} + \frac{-\mu_{21} - \mu_{22} + 1}{2(e^{\beta(-f_{2}(c) + f_{2}(d))} + 1)} &
-\frac{\mu_{11}}{2} - \frac{\mu_{21}}{2} + 1
- \frac{-\mu_{21} - \mu_{22} + 1}{2(e^{\beta(-f_{2}(c) + f_{2}(d))} + 1)}
- \frac{-\mu_{11} - \mu_{12} + 1}{2(e^{\beta(-f_{1}(b) + f_{1}(d))} + 1)}
\end{bmatrix}$}
\label{eqn:2x2_Introspection_Mutation_Transition}
\end{equation}

\subsection*{Introspective Imitation Dynamics with mutation}
Let \(\eta_{i,j}(\mathbf{x}, \mathbf{y}) = 1 + e^{\beta(f_i(\mathbf{x}) - f_j(\mathbf{y}))}\).
\begin{equation}
\resizebox{\textwidth}{!}{$
\begin{bmatrix}
-\frac{\mu_{12}}{2} - \frac{\mu_{22}}{2} + 1 &
\frac{\mu_{22}}{2} & \frac{\mu_{12}}{2} & 0\\[1em]
\frac{\mu_{21}}{2}
+ \frac{(\epsilon f_{1}(b) + 1)(-\mu_{21} - \mu_{22} + 1)}
       {2\eta_{2,2}(\mathbf{b},\mathbf{a})(\epsilon f_{1}(b) + \epsilon f_{2}(b) + 2)}
&
-\frac{\mu_{12}}{2} - \frac{\mu_{21}}{2}
- \frac{(\epsilon f_{1}(b) + 1)(-\mu_{21} - \mu_{22} + 1)}
       {2\eta_{2,2}(\mathbf{b},\mathbf{a})(\epsilon f_{1}(b) + \epsilon f_{2}(b) + 2)}
- \frac{(\epsilon f_{2}(b) + 1)(-\mu_{11} - \mu_{12} + 1)}
       {2\eta_{1,1}(\mathbf{b},\mathbf{d})(\epsilon f_{1}(b) + \epsilon f_{2}(b) + 2)}
+ 1 & 0 &
\frac{\mu_{12}}{2}
+ \frac{(\epsilon f_{2}(b) + 1)(-\mu_{11} - \mu_{12} + 1)}
       {2\eta_{1,1}(\mathbf{b},\mathbf{d})(\epsilon f_{1}(b) + \epsilon f_{2}(b) + 2)}\\[1em]
\frac{\mu_{11}}{2}
+ \frac{(\epsilon f_{2}(c) + 1)(-\mu_{11} - \mu_{12} + 1)}
       {2\eta_{1,1}(\mathbf{c},\mathbf{a})(\epsilon f_{1}(c) + \epsilon f_{2}(c) + 2)}
& 0 &
-\frac{\mu_{11}}{2} - \frac{\mu_{22}}{2}
- \frac{(\epsilon f_{1}(c) + 1)(-\mu_{21} - \mu_{22} + 1)}
       {2\eta_{2,2}(\mathbf{c},\mathbf{d})(\epsilon f_{1}(c) + \epsilon f_{2}(c) + 2)}
- \frac{(\epsilon f_{2}(c) + 1)(-\mu_{11} - \mu_{12} + 1)}
       {2\eta_{1,1}(\mathbf{c},\mathbf{a})(\epsilon f_{1}(c) + \epsilon f_{2}(c) + 2)}
+ 1 &
\frac{\mu_{22}}{2}
+ \frac{(\epsilon f_{1}(c) + 1)(-\mu_{21} - \mu_{22} + 1)}
       {2\eta_{2,2}(\mathbf{c},\mathbf{d})(\epsilon f_{1}(c) + \epsilon f_{2}(c) + 2)}\\[1em]
0 & \frac{\mu_{11}}{2} & \frac{\mu_{21}}{2} &
-\frac{\mu_{11}}{2} - \frac{\mu_{21}}{2} + 1
\end{bmatrix}$}
\label{eqn:2x2_IID_Mutation_Transition}
\end{equation}

\section{Mutation}


\subsection{The Moran Process}
In the case of the Moran process, our transition matrix becomes:

\begin{large}
    \begin{equation}
        T^{(m)}_\textbf{ab} = 
        \begin{cases}
            (\frac{\sum_{a_i = b_{I(\textbf{a}, \textbf{b})}}{f(a_i)}}{N\sum_{a_j}f(a_j)})(1-\sum_{k}\mu_{I(\textbf{a}, \textbf{b}),k}) + (\frac{\mu_{i,\textbf{b}_I(\textbf{a}, \textbf{b})}}{N}) & \text{if }\textbf{b} \in \text{Neb(\textbf{a})}\text{, differing at position }I(\textbf{a}, \textbf{b})\\
            0 & \text{if }\textbf{b} \notin \text{Neb(\textbf{a}) and \textbf{a}}\neq \textbf{b}\\
            1 - \sum_{\textbf{b} \in S \setminus \text{\{\textbf{a}\}}}T_{\textbf{ab}} & \text{if }\textbf{a}=\textbf{b}
        \end{cases}
    \label{eqn:Moran_Probability_Function_mutation}
    \end{equation}
\end{large}

The explicit transition matrix for the two-player case is given in
equation~(\ref{eqn:2x2_Moran_Mutation_Transition}) in the supplementary
material. This yields an ergodic Markov chain whose steady state is:
%TODO: fill in once computation on siren completes


\subsection{Fermi Imitation Dynamics}
When we introduce mutation to Fermi Imitation Dynamics, we obtain the following
transition matrix:

\begin{large}
    \begin{equation}
        T^{(m)}_\textbf{ab}  = 
        \begin{cases}
            (\frac{1}{N(N-1)}\sum_{a_j=b_{I(\textbf{a}, \textbf{b})}}\phi(\pi_{\text{I(\textbf{a}, \textbf{b})}}(a) - \pi_{j}(\textbf{a})))(1-\sum_{k}\mu_{I(\textbf{a}, \textbf{b}),k}) + (\frac{\mu_{i,\textbf{b}_I(\textbf{a}, \textbf{b})}}{N}) & \text{if }\textbf{b} \in \text{Neb(\textbf{a})}\\
            0 & \text{if }\textbf{b} \notin \text{Neb(\textbf{a}) and \textbf{a}}\neq \textbf{b}\\
            1 - \sum_{\textbf{b} \in S \setminus \text{\{\textbf{a}\}}}T_{\textbf{ab}} & \text{if }\textbf{a}=\textbf{b}
        \end{cases}
    \label{eqn:imitation_prob_mutation}
    \end{equation}
\end{large}


The explicit transition matrix for the two-player case is given in
equation~(\ref{eqn:2x2_Fermi_Mutation_Transition}) in the supplementary
material. The resulting steady state is:
%TODO: fill in once computation on siren completes

\subsection{Aspiration Dynamics}
As aspiration and introspection dynamics do not form absorbing Markov chains,
we do not need mutation to be applied in order to form an ergodic chain.
However, we will consider these cases for completeness, and thus we obtain the
following transition matrix:
\[
T^{(m)}_{\textbf{ab}} =
\begin{cases}
(\dfrac{1}{N} \, \phi (\pi_{I(\textbf{a}, \textbf{b})} - A_i))(1-\sum_{k}\mu_{I(\textbf{a}, \textbf{b}),k}) + (\frac{\mu_{i,\textbf{b}_I(\textbf{a}, \textbf{b})}}{N})
    & \text{if } \textbf{b} \in \mathrm{Neb}(\textbf{a}) \text{ and } j = I(\textbf{a}, \textbf{b}),\\[1.2em]
0
    & \text{if } \textbf{b} \notin \mathrm{Neb}(\textbf{a}) \text{and \textbf{a}} \neq \textbf{b},\\[0.8em]
1 - \sum_{\textbf{b} \in S \setminus \text{\{\textbf{a}\}}}T_{\textbf{ab}} & \text{if }\textbf{a}=\textbf{b}
\end{cases}
\]

The explicit transition matrix for the two-player case is given in
equation~(\ref{eqn:2x2_Aspiration_Mutation_Transition}) in the supplementary
material.


\subsection{Introspection Dynamics}
As introspection dynamics does not form an absorbing Markov chain, we do not
need mutation to form an ergodic chain.
However, we will consider the following transition matrix for the case with
mutation:

\[
T^{(m)}_{\textbf{ab}} =
\begin{cases}
(\dfrac{1}{N(m_j - 1)} \, \phi (\Delta \pi_{I(\textbf{a}, \textbf{b})}))(1-\sum_{k}\mu_{I(\textbf{a}, \textbf{b}),k}) + (\frac{\mu_{i,\textbf{b}_I(\textbf{a}, \textbf{b})}}{N})
    & \text{if } \textbf{b} \in \mathrm{Neb}(\textbf{a}) \text{ and } j = I(\textbf{a}, \textbf{b}),\\[1.2em]
0
    & \text{if } \textbf{b} \notin \mathrm{Neb}(\textbf{a}) \text{and \textbf{a}} \neq \textbf{b},\\[0.8em]
1 - \sum_{\textbf{b} \in S \setminus \text{\{\textbf{a}\}}}T_{\textbf{ab}} & \text{if }\textbf{a}=\textbf{b}
\end{cases}
\]

The explicit transition matrix for the two-player case is given in
equation~(\ref{eqn:2x2_Introspection_Mutation_Transition}) in the supplementary
material.

\subsection{Introspective Imitation Dynamics}
Finally, for the case of introspective imitation dynamics, we get the following
transition matrix with mutation:

\begin{equation}
T^{(m)}_{\textbf{ab}} =  \begin{cases}
    (\frac{1}{N}\frac{\sum_{a_{j} = b_{I(\textbf{a}, \textbf{b})}}f_j(\textbf{a})}{\sum_{k}f_k(\textbf{a})}\phi(\Delta(f_{I(\textbf{a}, \textbf{b})})))(1-\sum_{k}\mu_{I(\textbf{a}, \textbf{b}),k}) + (\frac{\mu_{i,\textbf{b}_I(\textbf{a}, \textbf{b})}}{N}) & \text{if \textbf{b}}\in \text{Neb(\textbf{a})}\\
    0 & \text{if \textbf{b}}\notin \text{Neb(\textbf{a}) and \textbf{a}} \neq \textbf{b}\\
    1 - \sum_{\textbf{b} \in S \setminus \text{\{\textbf{a}\}}}T_{\textbf{ab}} & \text{if }\textbf{a}=\textbf{b}
\end{cases}
\end{equation}

The explicit transition matrix for the two-player case is given in
equation~(\ref{eqn:2x2_IID_Mutation_Transition}) in the supplementary
material. The resulting steady state is:
%TODO: fill in once computation on siren completes

\subsection{Extended parameter sweeps}

\subsubsection*{Fermi imitation dynamics}
\begin{figure}[!h]
\centering

\begin{subfigure}{0.32\textwidth}
    \includegraphics[width=\linewidth]{../Figures/alpha_range_violinplots/Fermi_N_eq_2/main.pdf}
    \caption{$N=2$}
\end{subfigure}
\hfill
\begin{subfigure}{0.32\textwidth}
    \includegraphics[width=\linewidth]{../Figures/alpha_range_violinplots/Fermi_N_eq_3/main.pdf}
    \caption{$N=3$}
\end{subfigure}
\hfill
\begin{subfigure}{0.32\textwidth}
    \includegraphics[width=\linewidth]{../Figures/alpha_range_violinplots/Fermi_N_eq_4/main.pdf}
    \caption{$N=4$}
\end{subfigure}

\vspace{0.3cm}

\begin{subfigure}{0.32\textwidth}
    \includegraphics[width=\linewidth]{../Figures/alpha_range_violinplots/Fermi_N_eq_5/main.pdf}
    \caption{$N=5$}
\end{subfigure}
\hfill
\begin{subfigure}{0.32\textwidth}
    \includegraphics[width=\linewidth]{../Figures/alpha_range_violinplots/Fermi_N_eq_6/main.pdf}
    \caption{$N=6$}
\end{subfigure}
\hfill
\begin{subfigure}{0.32\textwidth}
    \includegraphics[width=\linewidth]{../Figures/alpha_range_violinplots/Fermi_N_eq_7/main.pdf}
    \caption{$N=7$}
\end{subfigure}

\vspace{0.3cm}

\begin{subfigure}{0.32\textwidth}
    \includegraphics[width=\linewidth]{../Figures/alpha_range_violinplots/Fermi_N_eq_8/main.pdf}
    \caption{$N=8$}
\end{subfigure}

\caption{Fermi imitation dynamics across population sizes $N$.}
\label{fig:fermi_full_grid}
\end{figure}


\subsubsection*{Introspective Imitation Dynamics}
\begin{figure}[!h]
\centering

\begin{subfigure}{0.32\textwidth}
    \includegraphics[width=\linewidth]{../Figures/alpha_range_violinplots/imispection_N_eq_2/main.pdf}
    \caption{$N=2$}
\end{subfigure}
\hfill
\begin{subfigure}{0.32\textwidth}
    \includegraphics[width=\linewidth]{../Figures/alpha_range_violinplots/imispection_N_eq_3/main.pdf}
    \caption{$N=3$}
\end{subfigure}
\hfill
\begin{subfigure}{0.32\textwidth}
    \includegraphics[width=\linewidth]{../Figures/alpha_range_violinplots/imispection_N_eq_4/main.pdf}
    \caption{$N=4$}
\end{subfigure}

\vspace{0.3cm}

\begin{subfigure}{0.32\textwidth}
    \includegraphics[width=\linewidth]{../Figures/alpha_range_violinplots/imispection_N_eq_5/main.pdf}
    \caption{$N=5$}
\end{subfigure}
\hfill
\begin{subfigure}{0.32\textwidth}
    \includegraphics[width=\linewidth]{../Figures/alpha_range_violinplots/imispection_N_eq_6/main.pdf}
    \caption{$N=6$}
\end{subfigure}
\hfill
\begin{subfigure}{0.32\textwidth}
    \includegraphics[width=\linewidth]{../Figures/alpha_range_violinplots/imispection_N_eq_7/main.pdf}
    \caption{$N=7$}
\end{subfigure}

\vspace{0.3cm}

\begin{subfigure}{0.32\textwidth}
    \includegraphics[width=\linewidth]{../Figures/alpha_range_violinplots/imispection_N_eq_8/main.pdf}
    \caption{$N=8$}
\end{subfigure}

\caption{Introspective imitation dynamics across population sizes $N$.}
\label{fig:introspection_grid}
\end{figure}


\subsubsection*{Moran process}
\begin{figure}[!h]
\centering

\begin{subfigure}{0.32\textwidth}
    \includegraphics[width=\linewidth]{../Figures/alpha_range_violinplots/Moran_N_eq_2/main.pdf}
    \caption{$N=2$}
\end{subfigure}
\hfill
\begin{subfigure}{0.32\textwidth}
    \includegraphics[width=\linewidth]{../Figures/alpha_range_violinplots/Moran_N_eq_3/main.pdf}
    \caption{$N=3$}
\end{subfigure}
\hfill
\begin{subfigure}{0.32\textwidth}
    \includegraphics[width=\linewidth]{../Figures/alpha_range_violinplots/Moran_N_eq_4/main.pdf}
    \caption{$N=4$}
\end{subfigure}

\vspace{0.3cm}

\begin{subfigure}{0.32\textwidth}
    \includegraphics[width=\linewidth]{../Figures/alpha_range_violinplots/Moran_N_eq_5/main.pdf}
    \caption{$N=5$}
\end{subfigure}
\hfill
\begin{subfigure}{0.32\textwidth}
    \includegraphics[width=\linewidth]{../Figures/alpha_range_violinplots/Moran_N_eq_6/main.pdf}
    \caption{$N=6$}
\end{subfigure}
\hfill
\begin{subfigure}{0.32\textwidth}
    \includegraphics[width=\linewidth]{../Figures/alpha_range_violinplots/Moran_N_eq_7/main.pdf}
    \caption{$N=7$}
\end{subfigure}

\vspace{0.3cm}

\begin{subfigure}{0.32\textwidth}
    \includegraphics[width=\linewidth]{../Figures/alpha_range_violinplots/Moran_N_eq_8/main.pdf}
    \caption{$N=8$}
\end{subfigure}

\caption{Moran process across population sizes $N$.}
\label{fig:moran_full_grid}
\end{figure}


\subsubsection*{Introspection dynamics}
\begin{figure}[!h]
\centering

\begin{subfigure}{0.32\textwidth}
    \includegraphics[width=\linewidth]{../Figures/alpha_range_violinplots/introspection_N_eq_2/main.pdf}
    \caption{$N=2$}
\end{subfigure}
\hfill
\begin{subfigure}{0.32\textwidth}
    \includegraphics[width=\linewidth]{../Figures/alpha_range_violinplots/introspection_N_eq_3/main.pdf}
    \caption{$N=3$}
\end{subfigure}
\hfill
\begin{subfigure}{0.32\textwidth}
    \includegraphics[width=\linewidth]{../Figures/alpha_range_violinplots/introspection_N_eq_4/main.pdf}
    \caption{$N=4$}
\end{subfigure}

\vspace{0.3cm}

\begin{subfigure}{0.32\textwidth}
    \includegraphics[width=\linewidth]{../Figures/alpha_range_violinplots/introspection_N_eq_5/main.pdf}
    \caption{$N=5$}
\end{subfigure}
\hfill
\begin{subfigure}{0.32\textwidth}
    \includegraphics[width=\linewidth]{../Figures/alpha_range_violinplots/introspection_N_eq_6/main.pdf}
    \caption{$N=6$}
\end{subfigure}
\hfill
\begin{subfigure}{0.32\textwidth}
    \includegraphics[width=\linewidth]{../Figures/alpha_range_violinplots/introspection_N_eq_7/main.pdf}
    \caption{$N=7$}
\end{subfigure}

\vspace{0.3cm}

\begin{subfigure}{0.32\textwidth}
    \includegraphics[width=\linewidth]{../Figures/alpha_range_violinplots/introspection_N_eq_8/main.pdf}
    \caption{$N=8$}
\end{subfigure}

\caption{Introspection dynamics across population sizes $N$.}
\label{fig:introspection_structure}
\end{figure}
\end{document}